% ===================================================================== PREAMBULE COMMUN — Carnet d'entraînement Fiche 9
\documentclass[11pt,a4paper]{article}
\usepackage[utf8]{inputenc}
\usepackage[T1]{fontenc}
\usepackage{lmodern}
\usepackage[french]{babel}
\usepackage{amsmath,amssymb,mathtools}
\usepackage{icomma}
\usepackage{siunitx}
\sisetup{output-decimal-marker={,},per-mode=symbol}
\DeclareSIUnit{\molar}{M}
\DeclareSIUnit{\Molar}{mol\,L^{-1}}
\usepackage[version=4]{mhchem}
\usepackage{xcolor}
\usepackage{colortbl}
\usepackage{tikz}
\usepackage{pgfplots}
\usepackage[most]{tcolorbox}
\usepackage{enumitem}
\usepackage{array}
\usepackage{booktabs}
\usepackage{multicol}
\usepackage{fancyhdr}
\usepackage[hidelinks]{hyperref}
\usetikzlibrary{arrows.meta,positioning,calc,decorations.pathmorphing,
  decorations.markings,angles,quotes,patterns,backgrounds,
  shapes.geometric,shapes.misc,fadings,intersections,fit}
\pgfplotsset{compat=1.18}
\usepackage[a4paper,margin=2.2cm,headheight=26pt,headsep=10pt]{geometry}

% ---------- Palette ----------
\definecolor{primary}{RGB}{150,98,162}
\definecolor{primarydark}{RGB}{92,52,110}
\definecolor{defcol}{RGB}{0,114,178}
\definecolor{thmcol}{RGB}{0,140,120}
\definecolor{formulacol}{RGB}{198,93,9}
\definecolor{remarkcol}{RGB}{120,94,200}
\definecolor{attncol}{RGB}{200,40,55}
\definecolor{excol}{RGB}{0,135,90}
\definecolor{methcol}{RGB}{90,90,100}
\definecolor{cationcol}{RGB}{0,90,190}
\definecolor{anioncol}{RGB}{200,60,55}
\definecolor{cristalcol}{RGB}{110,105,120}
\definecolor{eaucol}{RGB}{120,180,220}
\definecolor{solcol}{RGB}{0,135,90}
\definecolor{kpscol}{RGB}{198,93,9}
\definecolor{qcol}{RGB}{120,94,200}
\definecolor{precipcol}{RGB}{110,70,40}
\definecolor{exercol}{RGB}{150,98,162}
\definecolor{corrcol}{RGB}{0,120,100}

% ---------- Macros ----------
\newcommand{\Kps}{K_{\mathrm{ps}}}
\newcommand{\Ce}[1]{\left[\,\ce{#1}\,\right]}

% ---------- Style des boîtes ----------
\tcbset{boxbase/.style={enhanced,breakable,boxrule=0.9pt,arc=2.5pt,
   left=3mm,right=3mm,top=2.3mm,bottom=2.3mm,
   fonttitle=\bfseries,coltitle=white,toptitle=0.6mm,bottomtitle=0.6mm}}

\newtcolorbox[auto counter]{definition}[1][]{%
  boxbase,colback=defcol!5,colframe=defcol,
  title={\textsc{Définition}~\thetcbcounter\ \ifx\relax#1\relax\relax\else\textmd{--- \itshape #1}\fi}}
\newtcolorbox[auto counter]{theoreme}[1][]{%
  boxbase,colback=thmcol!6,colframe=thmcol,
  title={\textsc{Loi / Propriété}~\thetcbcounter\ \ifx\relax#1\relax\relax\else\textmd{--- \itshape #1}\fi}}
\newtcolorbox[auto counter]{formule}[1][]{%
  boxbase,colback=formulacol!6,colframe=formulacol,
  title={\textsc{Formule}~\thetcbcounter\ \ifx\relax#1\relax\relax\else\textmd{--- \itshape #1}\fi}}
\newtcolorbox{remarque}[1][]{%
  boxbase,colback=remarkcol!6,colframe=remarkcol,
  title={\textsc{Remarque}\ifx\relax#1\relax\relax\else\textmd{ : \itshape #1}\fi}}
\newtcolorbox{attention}[1][]{%
  boxbase,colback=attncol!6,colframe=attncol,
  title={\textsc{Attention}\ifx\relax#1\relax\relax\else\textmd{ : \itshape #1}\fi}}
\newtcolorbox{exemple}[1][]{%
  boxbase,colback=excol!5,colframe=excol,
  title={\textsc{Exemple}\ifx\relax#1\relax\relax\else\textmd{ : \itshape #1}\fi}}
\newtcolorbox{methode}[1][]{%
  boxbase,colback=methcol!7,colframe=methcol,sharp corners,
  title={\textsc{Méthode}\ifx\relax#1\relax\relax\else\textmd{ : \itshape #1}\fi}}
\newtcolorbox{astuce}[1][]{%
  boxbase,colback=primary!7,colframe=primary,
  title={\textsc{Astuce}\ifx\relax#1\relax\relax\else\textmd{ : \itshape #1}\fi}}

\newtcolorbox[auto counter]{exercice}[1][]{%
  boxbase,colback=exercol!4,colframe=exercol,
  title={\textsc{Exercice}~\thetcbcounter\ \ifx\relax#1\relax\relax\else\textmd{--- \itshape #1}\fi}}
\newtcolorbox[auto counter]{correction}[1][]{%
  boxbase,colback=corrcol!4,colframe=corrcol,
  title={\textsc{Corrigé}~\thetcbcounter\ \ifx\relax#1\relax\relax\else\textmd{--- \itshape #1}\fi}}

% ---------- Environnement « où : » ----------
\newenvironment{ou}{%
  \par\smallskip\noindent\textbf{où :}\nopagebreak
  \begin{itemize}[leftmargin=1.8em,topsep=2pt,itemsep=1.5pt,parsep=0pt]}%
  {\end{itemize}}

% ---------- Styles TikZ ----------
\tikzset{
  cat/.style={circle,fill=cationcol,draw=cationcol!50!black,inner sep=0pt,
              minimum size=5mm,font=\bfseries\scriptsize,text=white},
  an/.style={circle,fill=anioncol,draw=anioncol!50!black,inner sep=0pt,
             minimum size=5mm,font=\bfseries\scriptsize,text=white},
  crx/.style={circle,fill=cristalcol,draw=black!50,inner sep=0pt,
              minimum size=5mm,font=\bfseries\scriptsize,text=white},
  ptn/.style={circle,fill=black,inner sep=1.1pt}}

% ---------- En-têtes ----------
\pagestyle{fancy}
\fancyhf{}
\fancyhead[L]{\small\textcolor{primarydark}{\textbf{Fiche 9} — Solubilité et produit de solubilité}}
\fancyhead[R]{\small\textcolor{primarydark}{\textsc{Carnet d'entraînement}}}
\fancyfoot[C]{\small\thepage}
\renewcommand{\headrulewidth}{0.8pt}
\renewcommand{\headrule}{\hbox to\headwidth{\color{primary}\leaders\hrule height \headrulewidth\hfill}}
\usepackage{titlesec}
\titleformat{\section}{\Large\bfseries\color{primarydark}}{\thesection}{0.6em}{}
\titleformat{\subsection}{\large\bfseries\color{primary}}{\thesubsection}{0.6em}{}

% ---------- Bandeau de titre ----------
% #1 = thème/étiquette   #2 = titre principal   #3 = sous-titre
\newcommand{\bandeau}[3]{%
\begin{center}
\begin{tikzpicture}
\node[rectangle,rounded corners=7pt,fill=primary,text=white,
      minimum width=15.6cm,minimum height=1.95cm,align=center,inner sep=8pt]
  {{\large\bfseries #1}\\[3pt]{\Large\bfseries #2}\\[3pt]{\small #3}};
\end{tikzpicture}
\end{center}
\vspace{2mm}}
\begin{document}
\bandeau{Thème 4 — Effet d'ion commun}%
{Corrigé — Niveau Moyen}%
{Réponses détaillées et justifications}

\begin{correction}[ion commun anionique — \ce{CaF2}]
$\ce{CaF2(s) <=> Ca^{2+} + 2 F-}$, $\Kps=[\ce{Ca^2+}][\ce{F-}]^{2}$.
\begin{enumerate}[label=\alph*),leftmargin=2em,topsep=2pt,itemsep=2pt]
  \item Ion commun : \ce{F-}. Comme $[\ce{F-}]\approx\qty{0{,}10}{\molar}$ (imposé) et $[\ce{Ca^2+}]=s'$ :
        $\Kps\approx s'\,[\ce{F-}]^{2}$, donc $s'=\dfrac{\Kps}{[\ce{F-}]^{2}}$.
  \item $s'=\dfrac{\num{3{,}9e-11}}{(0{,}10)^{2}}=\dfrac{\num{3{,}9e-11}}{\num{1{,}0e-2}}=\boxed{\qty{3{,}9e-9}{\molar}}$.
  \item $\dfrac{s_0}{s'}=\dfrac{\num{2{,}1e-4}}{\num{3{,}9e-9}}\approx\num{5{,}4e4}$ : la solubilité est
        divisée par environ \num{54000}.
\end{enumerate}
\end{correction}

\begin{correction}[ion commun cationique — \ce{BaSO4}]
\begin{enumerate}[label=\alph*),leftmargin=2em,topsep=2pt,itemsep=2pt]
  \item $\Kps=[\ce{Ba^2+}][\ce{SO4^2-}]$ ; avec $[\ce{Ba^2+}]$ imposé et $[\ce{SO4^2-}]=s'$ :
        $s'=\dfrac{\Kps}{[\ce{Ba^2+}]}$.
  \item $s'=\dfrac{\num{1{,}1e-10}}{\num{1{,}0e-2}}=\boxed{\qty{1{,}1e-8}{\molar}}$.
        Or $s_0=\sqrt{\Kps}=\sqrt{\num{1{,}1e-10}}\approx\qty{1{,}0e-5}{\molar}$ : la solubilité est
        divisée par $\sim\num{1000}$. L'ion commun peut donc être le cation aussi bien que l'anion.
\end{enumerate}
\end{correction}

\begin{correction}[ion commun \ce{OH-} — \ce{Mg(OH)2}]
$\ce{Mg(OH)2(s) <=> Mg^{2+} + 2 OH-}$, type $1{:}2$.
\begin{enumerate}[label=\alph*),leftmargin=2em,topsep=2pt,itemsep=2pt]
  \item $\Kps=[\ce{Mg^2+}][\ce{OH-}]^{2}\approx s'\,[\ce{OH-}]^{2}$, donc $s'=\dfrac{\Kps}{[\ce{OH-}]^{2}}$.
  \item $s'=\dfrac{\num{9{,}0e-12}}{(\num{1{,}0e-2})^{2}}=\dfrac{\num{9{,}0e-12}}{\num{1{,}0e-4}}
        =\boxed{\qty{9{,}0e-8}{\molar}}$.
\end{enumerate}
\emph{(À rapprocher du Thème 5 : ajouter des \ce{OH-}, c'est augmenter le pH, ce qui fait précipiter les
hydroxydes.)}
\end{correction}

\begin{correction}[dépendance en $C$]
$s'=\Kps/[\ce{Cl-}]$ :
\begin{enumerate}[label=\alph*),leftmargin=2em,topsep=2pt,itemsep=2pt]
  \item $s'=\dfrac{\num{1{,}8e-10}}{\num{2{,}0e-2}}=\qty{9{,}0e-9}{\molar}$.
  \item $s'=\dfrac{\num{1{,}8e-10}}{0{,}10}=\qty{1{,}8e-9}{\molar}$.
  \item En multipliant $[\ce{Cl-}]$ par 5 (de \num{0{,}02} à \num{0{,}10}), $s'$ est \textbf{divisée par
        5} (de \num{9{,}0e-9} à \num{1{,}8e-9}) : $s'$ est \emph{inversement proportionnelle} à
        l'ion commun (pour un sel $1{:}1$).
\end{enumerate}
\end{correction}

\begin{correction}[le $\Kps$ ne bouge pas]
\begin{enumerate}[label=\alph*),leftmargin=2em,topsep=2pt,itemsep=2pt]
  \item $[\ce{Ca^2+}]=s'=\qty{3{,}9e-9}{\molar}$ ; $[\ce{F-}]\approx\qty{0{,}10}{\molar}$ (l'apport de la
        dissolution, $2s'\approx\num{7{,}8e-9}$, est négligeable devant \num{0,10}).
  \item $[\ce{Ca^2+}][\ce{F-}]^{2}=(\num{3{,}9e-9})(0{,}10)^{2}=(\num{3{,}9e-9})(\num{1{,}0e-2})
        =\num{3{,}9e-11}=\Kps$. \checkmark\;
        Le \textbf{$\Kps$ est inchangé} (constante à $T$ fixée) ; c'est la \textbf{solubilité} qui a
        chuté pour maintenir le produit constant malgré l'excès de \ce{F-}.
\end{enumerate}
\end{correction}

\end{document}
